% libre-compendium
%
% Copyright (C) 2026 Richard Qin
%
% This document is dual-licensed under either:
% - GNU Free Documentation License, Version 1.3 or any later version
%   published by the Free Software Foundation; with no Invariant Sections,
%   no Front-Cover Texts, and no Back-Cover Texts (see COPYING.DOC); or
% - Creative Commons Attribution-ShareAlike 4.0 International License
%   (see LICENSE-CC-BY-SA-4.0).
\section{随机变量及其分布}

\begin{exercise}[title={离散型分布律的建立},level=basic,points=6,tags={分布律,古典概型}]
(1) 袋中有编号为 $1,2,3,4,5$ 的 $5$ 只球，从中同时任取 $3$ 只，以 $X$ 表示取出 $3$ 只球中的最大号码，写出 $X$ 的分布律。

(2) 抛掷一颗均匀骰子两次，以 $Y$ 表示两次点数中的较小值，写出 $Y$ 的分布律。
\begin{answer}
(1) $P(X=3)=\frac{1}{10}$，$P(X=4)=\frac{3}{10}$，$P(X=5)=\frac{6}{10}$；
(2) $P(Y=k)=\frac{13-2k}{36}\,(k=1,2,3,4,5,6)$。
\end{answer}
\begin{solution}
(1) $X$ 的可能取值为 $3,4,5$，总取法为 $\binom{5}{3}=10$。最大号码为 $k$ 意味着必取 $k$ 号球，且另两球从 $1,\ldots,k-1$ 中选取，共有 $\binom{k-1}{2}$ 种取法，故
\[
  P(X=3)=\frac{\binom{2}{2}}{10}=\frac{1}{10},\qquad
  P(X=4)=\frac{\binom{3}{2}}{10}=\frac{3}{10},\qquad
  P(X=5)=\frac{\binom{4}{2}}{10}=\frac{6}{10}.
\]

(2) 有序样本空间含 $36$ 个等可能点。$Y=k$ 表示两数均不小于 $k$ 且至少一个等于 $k$，即 $(k,k)$ 以及 $(k,j),(j,k)\,(j=k+1,\ldots,6)$，共 $1+2(6-k)=13-2k$ 个样本点，故
\[
  P(Y=k)=\frac{13-2k}{36},\qquad k=1,2,3,4,5,6.
\]
核验 $\sum_{k=1}^{6}(13-2k)=11+9+7+5+3+1=36$，总概率为 $1$。
\end{solution}
\end{exercise}

\begin{exercise}[title={二项分布的基本计算},level=basic,points=6,tags={二项分布}]
一大楼装有 $5$ 台同类型的供水设备，各台设备是否被使用相互独立，任一时刻每台被使用的概率均为 $0.1$。求同一时刻：(1) 恰有 $2$ 台被使用的概率；(2) 至少有 $3$ 台被使用的概率；(3) 至少有 $1$ 台被使用的概率。
\begin{answer}
(1) $0.0729$；(2) $0.00856$；(3) $0.40951$。
\end{answer}
\begin{solution}
记同一时刻被使用的设备台数为 $X$，则 $X\sim b(5,0.1)$。
(1) $P(X=2)=\binom{5}{2}(0.1)^2(0.9)^3=10\times 0.01\times 0.729=0.0729$。
(2)
\[
\begin{aligned}
  P(X\ge 3)&=P(X=3)+P(X=4)+P(X=5)\\
  &=\binom{5}{3}(0.1)^3(0.9)^2+\binom{5}{4}(0.1)^4(0.9)+(0.1)^5\\
  &=0.0081+0.00045+0.00001=0.00856.
\end{aligned}
\]
(3) 利用对立事件：$P(X\ge 1)=1-P(X=0)=1-0.9^5=1-0.59049=0.40951$。
\end{solution}
\end{exercise}

\begin{exercise}[title={泊松分布与泊松近似},level=basic,points=6,tags={泊松分布,泊松定理}]
(1) 某电话总机每分钟收到呼唤的次数 $X\sim\pi(4)$，求某一分钟内呼唤次数大于 $3$ 的概率。

(2) 一保险公司承保了 $5000$ 份同年龄段一年期寿险，每人在一年内死亡的概率为 $0.0015$ 且相互独立。若每人死亡赔付 $3$ 万元，用泊松定理近似求公司赔付总额不超过 $30$ 万元的概率（可保留求和式并给出近似数值）。
\begin{answer}
(1) $1-\frac{71}{3}\e^{-4}\approx 0.5665$；(2) $\sum_{k=0}^{10}\frac{7.5^k\e^{-7.5}}{k!}\approx 0.8622$。
\end{answer}
\begin{solution}
(1) 利用对立事件，
\[
  P(X>3)=1-\sum_{k=0}^{3}\frac{4^k\e^{-4}}{k!}
  =1-\e^{-4}\left(1+4+8+\frac{32}{3}\right)
  =1-\frac{71}{3}\e^{-4}\approx 0.5665.
\]

(2) 记一年内死亡人数为 $Y\sim b(5000,0.0015)$，赔付总额不超过 $30$ 万元等价于 $Y\le 10$。取 $\lambda=np=5000\times 0.0015=7.5$，由泊松近似得
\[
  P(Y\le 10)\approx\sum_{k=0}^{10}\frac{7.5^k\e^{-7.5}}{k!}\approx 0.8622.
\]
\end{solution}
\end{exercise}

\begin{exercise}[title={连续型分布函数与概率密度},level=basic,points=6,tags={分布函数,概率密度}]
设连续型随机变量 $X$ 的分布函数为
\[
  F(x)=
  \begin{cases}
    0,& x<1,\\
    \ln x,& 1\le x<\e,\\
    1,& x\ge\e.
  \end{cases}
\]
(1) 求 $P(X<2)$、$P(0<X<3)$ 与 $P(2<X<5/2)$；
(2) 求 $X$ 的概率密度 $f(x)$。
\begin{answer}
(1) $\ln 2$，$1$，$\ln(5/4)$；(2) $f(x)=1/x\,(1<x<\e)$，其余为 $0$。
\end{answer}
\begin{solution}
(1) 因 $X$ 为连续型，区间端点开闭不影响概率：
\[
  P(X<2)=F(2)=\ln 2,\qquad
  P(0<X<3)=F(3)-F(0)=1-0=1,
\]
\[
  P(2<X<5/2)=F(5/2)-F(2)=\ln\frac{5}{2}-\ln 2=\ln\frac{5}{4}.
\]

(2) 对 $F(x)$ 分段求导，得
\[
  f(x)=
  \begin{cases}
    \dfrac{1}{x},& 1<x<\e,\\
    0,& \text{其他}.
  \end{cases}
\]
\end{solution}
\end{exercise}

\begin{exercise}[title={几何分布与偶次成功概率},level=standard,points=8,tags={几何分布,级数求和}]
一篮球运动员每次投篮命中率为 $p=0.45$，各次投篮相互独立。以 $X$ 表示他首次投中时累计已投篮的次数。
(1) 写出 $X$ 的分布律；
(2) 求 $X$ 取偶数的概率。
\begin{answer}
(1) $P(X=k)=0.45\times 0.55^{k-1}\,(k=1,2,\ldots)$；(2) $\frac{11}{31}\approx 0.3548$。
\end{answer}
\begin{solution}
(1) 显然 $X\sim G(0.45)$，分布律为
\[
  P(X=k)=0.45\times 0.55^{k-1},\qquad k=1,2,\ldots
\]

(2) 对所有正偶数求和，利用等比级数公式（记 $q=1-p=0.55$）：
\[
  P(X\text{ 为偶数})=\sum_{m=1}^{\infty}P(X=2m)
  =\sum_{m=1}^{\infty}pq^{2m-1}
  =\frac{pq}{1-q^2}
  =\frac{q}{1+q}
  =\frac{0.55}{1.55}=\frac{11}{31}.
\]
\end{solution}
\end{exercise}

\begin{exercise}[title={麦克斯韦速度分布的归一化},level=standard,points=8,tags={概率密度,广义积分}]
气体分子运动速率 $X$ 服从麦克斯韦（Maxwell）分布，其概率密度为
\[
  f(x)=
  \begin{cases}
    Ax^2\e^{-bx^2},& x>0,\\
    0,& x\le 0,
  \end{cases}
\]
其中 $b>0$ 为常数，求归一化常数 $A$。
\begin{answer}
$A=\frac{4b^{3/2}}{\sqrt{\pi}}$。
\end{answer}
\begin{solution}
由规范性 $A\int_{0}^{+\infty}x^2\e^{-bx^2}\dd x=1$。分部积分（令 $u=x,\,\dd v=x\e^{-bx^2}\dd x$，则 $v=-\frac{1}{2b}\e^{-bx^2}$）：
\[
  \int_{0}^{+\infty}x^2\e^{-bx^2}\dd x
  =\left[-\frac{x}{2b}\e^{-bx^2}\right]_{0}^{+\infty}
  +\frac{1}{2b}\int_{0}^{+\infty}\e^{-bx^2}\dd x.
\]
令 $t=\sqrt{2b}\,x$，由高斯积分 $\int_{0}^{+\infty}\e^{-t^2/2}\dd t=\sqrt{\pi/2}$ 得
\[
  \int_{0}^{+\infty}\e^{-bx^2}\dd x=\frac{1}{\sqrt{2b}}\sqrt{\frac{\pi}{2}}=\frac{\sqrt{\pi}}{2\sqrt{b}}.
\]
于是原积分等于 $\frac{\sqrt{\pi}}{4b^{3/2}}$，解得 $A=\frac{4b^{3/2}}{\sqrt{\pi}}$。
\end{solution}
\end{exercise}

\begin{exercise}[title={正态分布的区间概率与参数反求},level=standard,points=8,tags={正态分布,分位点}]
(1) 设 $X\sim N(3,2^2)$，求 $P(2<X<5)$、$P(\abs{X}>2)$，并确定常数 $d$ 使 $P(X>d)=0.9$。

(2) 某元件寿命 $T\sim N(160,\sigma^2)$，若要求 $P(120<T<200)\ge 0.80$，允许的标准差 $\sigma$ 最大为多少？（参考值：$\Phi(0.5)\approx 0.6915$，$\Phi(1)\approx 0.8413$，$\Phi(2.5)\approx 0.9938$，$z_{0.10}\approx 1.282$。）
\begin{answer}
(1) $0.5328$，$0.6977$，$d\approx 0.436$；(2) $\sigma\le 31.20$。
\end{answer}
\begin{solution}
(1) 由标准化公式（$\mu=3,\,\sigma=2$）：
\[
  P(2<X<5)=\Phi(1)-\Phi(-0.5)=\Phi(1)+\Phi(0.5)-1\approx 0.8413+0.6915-1=0.5328.
\]
对 $P(\abs{X}>2)=1-P(-2\le X\le 2)$，
\[
  P(-2\le X\le 2)=\Phi(-0.5)-\Phi(-2.5)=\Phi(2.5)-\Phi(0.5)\approx 0.9938-0.6915=0.3023,
\]
故 $P(\abs{X}>2)\approx 0.6977$。由 $P(X>d)=1-\Phi\bigl(\frac{d-3}{2}\bigr)=0.9$ 得 $\Phi\bigl(\frac{3-d}{2}\bigr)=0.9$，即 $\frac{3-d}{2}=z_{0.10}\approx 1.282$，解得 $d\approx 3-2.564=0.436$。

(2) $P(120<T<200)=\Phi\bigl(\frac{40}{\sigma}\bigr)-\Phi\bigl(-\frac{40}{\sigma}\bigr)=2\Phi\bigl(\frac{40}{\sigma}\bigr)-1\ge 0.80$，即 $\Phi\bigl(\frac{40}{\sigma}\bigr)\ge 0.90$。故 $\frac{40}{\sigma}\ge z_{0.10}\approx 1.282$，得 $\sigma\le\frac{40}{1.282}\approx 31.20$。
\end{solution}
\end{exercise}

\begin{exercise}[title={均匀分布与二次方程实根},level=standard,points=8,tags={均匀分布,综合建模}]
设随机变量 $K\sim U(0,5)$，求关于 $x$ 的一元二次方程 $4x^2+4Kx+K+2=0$ 有实根的概率。
\begin{answer}
$\frac{3}{5}$。
\end{answer}
\begin{solution}
方程有实根当且仅当判别式非负：
\[
  \Delta=(4K)^2-4\cdot 4\cdot(K+2)=16(K^2-K-2)=16(K-2)(K+1)\ge 0,
\]
即 $K\le -1$ 或 $K\ge 2$。结合 $K\sim U(0,5)$ 的有效区间 $(0,5)$，所求概率为
\[
  P(K\ge 2)=\int_{2}^{5}\frac{1}{5}\dd k=\frac{3}{5}.
\]
\end{solution}
\end{exercise}

\begin{exercise}[title={单调与非单调连续函数变换},level=standard,points=10,tags={随机变量函数}]
(1) 设 $X\sim U(0,1)$，求 $Y=-2\ln X$ 的概率密度；

(2) 设 $X\sim N(0,1)$，求 $Y=\e^X$（对数正态分布）与 $W=\abs{X}$ 的概率密度。
\begin{answer}
(1) $f_Y(y)=\frac{1}{2}\e^{-y/2}\,(y>0)$，其余为 $0$（即 $\operatorname{Exp}(2)$）；
(2) $f_Y(y)=\frac{1}{\sqrt{2\pi}\,y}\e^{-(\ln y)^2/2}\,(y>0)$，$f_W(w)=\sqrt{\frac{2}{\pi}}\e^{-w^2/2}\,(w>0)$，其余为 $0$。
\end{answer}
\begin{solution}
(1) $g(x)=-2\ln x$ 在 $(0,1)$ 上严格递减，值域为 $(0,+\infty)$，反函数 $h(y)=\e^{-y/2}$，$\abs{h'(y)}=\frac{1}{2}\e^{-y/2}$。由单调变换公式得
\[
  f_Y(y)=
  \begin{cases}
    \dfrac{1}{2}\e^{-y/2},& y>0,\\
    0,& y\le 0,
  \end{cases}
\]
即 $Y\sim\operatorname{Exp}(2)$。

(2) 对 $Y=\e^X$，值域为 $(0,+\infty)$，反函数 $h(y)=\ln y$，$h'(y)=1/y$，故对 $y>0$，
\[
  f_Y(y)=\varphi(\ln y)\cdot\frac{1}{y}=\frac{1}{\sqrt{2\pi}\,y}\e^{-(\ln y)^2/2},
\]
$y\le 0$ 时密度为 $0$。

对 $W=\abs{X}$，当 $w\le 0$ 时 $F_W(w)=0$；当 $w>0$ 时，
\[
  F_W(w)=P(-w\le X\le w)=2\Phi(w)-1.
\]
求导得
\[
  f_W(w)=
  \begin{cases}
    2\varphi(w)=\sqrt{\dfrac{2}{\pi}}\,\e^{-w^2/2},& w>0,\\
    0,& w\le 0.
  \end{cases}
\]
\end{solution}
\end{exercise}

\begin{exercise}[title={混合型随机变量的判别},level=advanced,points=8,tags={分布函数,混合分布}]
某人上班途经一处红绿灯，该路口有 $80\%$ 的时间亮红灯；遇红灯时等待时间（单位：$\mathrm{s}$）在 $[0,30]$ 上均匀分布，遇绿灯则无需等待。记 $X$ 为他的等待时间，求 $X$ 的分布函数 $F(x)$，并判断 $X$ 是离散型、连续型还是混合型随机变量。
\begin{answer}
$F(x)=0\,(x<0)$，$0.2+\frac{2}{75}x\,(0\le x<30)$，$1\,(x\ge 30)$；混合型。
\end{answer}
\begin{solution}
记 $A$ 为“遇绿灯”，$P(A)=0.2$，$P(\overline{A})=0.8$。当 $x<0$ 时 $F(x)=0$；当 $x\ge 30$ 时 $F(x)=1$；当 $0\le x<30$ 时，由全概率公式：
\[
  F(x)=P(A)P(X\le x\mid A)+P(\overline{A})P(X\le x\mid\overline{A})
  =0.2\times 1+0.8\times\frac{x}{30}=0.2+\frac{2x}{75}.
\]
综上，
\[
  F(x)=
  \begin{cases}
    0,& x<0,\\
    0.2+\dfrac{2}{75}x,& 0\le x<30,\\
    1,& x\ge 30.
  \end{cases}
\]
由于 $F(0)-F(0-0)=0.2>0$，$F$ 在 $x=0$ 处间断，故 $X$ 不是连续型；又 $X$ 在 $(0,30)$ 内连续取值，不是可列集合，故也不是离散型。因此 $X$ 是\textbf{混合型}随机变量。
\end{solution}
\end{exercise}

\begin{exercise}[title={概率密度的凸组合},level=advanced,points=8,tags={概率密度,证明}]
设 $f(x)$ 与 $g(x)$ 均为概率密度函数，$0\le\alpha\le 1$ 为常数。证明 $h(x)=\alpha f(x)+(1-\alpha)g(x)$ 也是概率密度函数，并说明其概率模型解释。
\begin{answer}
见解答。
\end{answer}
\begin{solution}
由 $f(x)\ge 0$、$g(x)\ge 0$ 及 $\alpha\ge 0$、$1-\alpha\ge 0$，立即得 $h(x)\ge 0$。由积分线性性及 $f,g$ 的规范性，
\[
  \int_{-\infty}^{+\infty}h(x)\dd x
  =\alpha\int_{-\infty}^{+\infty}f(x)\dd x+(1-\alpha)\int_{-\infty}^{+\infty}g(x)\dd x
  =\alpha\cdot 1+(1-\alpha)\cdot 1=1.
\]
由概率密度的判别条件可知 $h(x)$ 必为某个连续型随机变量的概率密度函数。在模型解释上，$h(x)$ 对应一个二元混合模型：先以概率 $\alpha$ 选择分布 $f$、以概率 $1-\alpha$ 选择分布 $g$，再从所选分布中抽取样本。
\end{solution}
\end{exercise}

\begin{exercise}[title={非单调三角函数变换},level=advanced,points=10,tags={随机变量函数,多分支}]
设随机变量 $X$ 具有概率密度
\[
  f_X(x)=
  \begin{cases}
    \dfrac{1}{\pi},& 0<x<\pi,\\
    0,& \text{其他}.
  \end{cases}
\]
求 $Y=\sin X$ 的概率密度 $f_Y(y)$。
\begin{answer}
$f_Y(y)=\frac{2}{\pi\sqrt{1-y^2}}\,(0<y<1)$，其余为 $0$。
\end{answer}
\begin{solution}
当 $x\in(0,\pi)$ 时，$y=\sin x\in(0,1]$，且在 $(0,\pi/2)$ 递增、在 $(\pi/2,\pi)$ 递减（不单调），采用分布函数法求解。当 $y\le 0$ 时 $F_Y(y)=0$；当 $y\ge 1$ 时 $F_Y(y)=1$。当 $0<y<1$ 时，方程 $\sin x=y$ 在 $(0,\pi)$ 内有两个根 $\arcsin y$ 与 $\pi-\arcsin y$，从而
\[
\begin{aligned}
  F_Y(y)&=P(\sin X\le y)
  =P(0<X\le\arcsin y)+P(\pi-\arcsin y\le X<\pi)\\
  &=\frac{\arcsin y}{\pi}+\frac{\arcsin y}{\pi}=\frac{2}{\pi}\arcsin y.
\end{aligned}
\]
对 $y\in(0,1)$ 求导得
\[
  f_Y(y)=
  \begin{cases}
    \dfrac{2}{\pi\sqrt{1-y^2}},& 0<y<1,\\
    0,& \text{其他}.
  \end{cases}
\]
\end{solution}
\end{exercise}
